problem:  
Let \(G = \mathrm{SU}(2) \times \mathrm{SU}(2)\) act on a compact, simply connected 7‑manifold \(M^7\) with cohomogeneity one and orbit diagram
\[
G \supset K_-, K_+ \supset H,
\]
where \(K_\pm/H \cong \mathbb{S}^2\), so that \(M\) belongs to the Grove–Ziller family of nonnegatively curved cohomogeneity‑one 7‑manifolds (specifically, certain \(\mathbb{S}^3\)-bundles over \(\mathbb{S}^4\)).  
Denote by \(\mathcal{M}_{G,\ge 0}(M)\) the moduli space of all \(G\)-invariant Riemannian metrics on \(M\) with nonnegative sectional curvature, modulo global scaling and \(G\)-equivariant diffeomorphism equivalence.

**Open Problem (Moduli Rigidity for Grove–Ziller 7‑Manifolds):**  
Is the moduli space \(\mathcal{M}_{G,\ge 0}(M)\) discrete, or does it contain a positive‑dimensional connected component?  
Equivalently, does there exist a smooth one‑parameter family of \(G\)-invariant, nonnegatively curved metrics \(\{g_s\}_{s\in[0,1]}\) on \(M\) that are pairwise inequivalent under any \(G\)-equivariant diffeomorphism and scalar multiple?

If such continuous families exist, what are their infinitesimal deformation directions in the space of invariant warped metrics \(dt^2 + g_t\)?  Do these deformations necessarily involve varying the warping functions near the singular orbits in a curvature‑preserving way, or can they arise from global Cheeger‑type deformations?

Why is it a "good" problem:  
This problem refines the earlier general question into a **specific and concrete case**—the best‑understood family of compact, simply connected cohomogeneity‑one manifolds supporting nonnegative curvature, yet with almost no information about their moduli space structure.  

1. **Conceptually Deep:**  
   It interrogates whether nonnegative curvature and maximal symmetry enforce rigidity of the metric structure or allow flexible deformations.  Even for these explicitly constructed \(\mathrm{SU}(2)\times\mathrm{SU}(2)\) actions, no result currently describes the dimension or topology of \(\mathcal{M}_{G,\ge 0}(M)\).  

2. **Bridge Across Fields:**  
   The problem unites Lie group actions (via \(G\supset K_\pm \supset H\)), global analysis of curvature inequalities (nonlinear PDEs defining \(\sec\ge 0\)), and moduli‑space theory reminiscent of those in Einstein or Ricci soliton geometry.  Any progress would likely require developing new analytic or topological tools to understand deformations respecting curvature bounds and symmetry constraints.  

3. **Simple but Profound:**  
   The formulation—“Is the moduli space discrete or continuous?”—is immediately understandable and geometrically natural, yet addresses an issue that remains completely unresolved even in this highly symmetric, accessible manifold class.  

4. **Forward‑Looking Impact:**  
   A rigid (discrete) moduli outcome would highlight an unexpected stability phenomenon in the landscape of nonnegatively curved metrics with symmetry, analogous to Mostow‑type rigidity.  A flexible (continuous) outcome would reveal uncharted moduli behavior and potentially new geometric flows preserving nonnegative curvature.  Either direction will fundamentally deepen our structural understanding of invariant curvature geometry.